% !TEX root = Main.tex
%

\section{Lower bound in the best of both worlds} %setting (Theorem~\ref{th:lowBOB}) }
\label{app:prooflowBOB}
%
\setcounter{scratchcounter}{\value{theorem}}\setcounter{theorem}{\the\numexpr\getrefnumber{th:lowBOB}-1}
\thi*
\setcounter{theorem}{\the\numexpr\value{scratchcounter}}
\begin{proof}%[Theorem~\ref{th:lowBOB}]\
	%    
As in Theorem 2, we exploit the fact that two environments with different best
arms can produce observations that are difficult to distinguish. Here, the
stochastic guarantee forces the learner to recommend one arm with very high
probability, while the adversarial environment requires a different
recommendation. We use a change-of-measure argument to quantify this conflict,
then show that it persists for a deterministic adversarial reward table.

Fix the reference gaps defining $\Gaps_4$, a learner, and a
horizon $\timeHorizon$ satisfying the hypothesis of the theorem.
Throughout the proof, the learner is run with this fixed horizon,
which it may know. The constructions below may depend on this
learner and on $\timeHorizon$.
Let $\pullsNumber_k(t)$ be the number of times arm $k$ has been
pulled by the end of round $t$.
	Let us consider a base problem, denoted \basePro{}, with $\nu_1,\ldots, \nu_\nArms$, 
	Bernoulli distributions with mean $\mean^{\basePro}_1,
	\ldots, \mean^{\basePro}_\nArms$ 
	such that $\mean^{\basePro}_1 \triangleq 1/2$, and for all $k\in[2,\nArms], 
	\mean^{\basePro}_k \triangleq  1/2-\gap_k$ 
	(the gaps specified in the claim of the theorem).
	The expectation and probability with respect to the
	learner and the samples of this problem are 
	$\Exp_{\basePro}$ and $\Pro_{\basePro}$.% with $\gap_k\geq 0$ for .
	
	With respect to this base problem, choose any maximizing index $i\in\arg\max_{j\in\setArms}\frac{j}{\gap_j}$.
 %denotes the rank of a suboptimal arm in the	base problem. 
\noindent\textbf{Choice of the initial phase length.}
We choose the length of the initial phase so as to maximize the resulting
lower bound on the learner's error. The stochastic and adversarial
constructions differ only during this phase. A shorter phase reduces
the phase-length contribution $\timeHorizon_i\gap_i^2/i$ to the
information-theoretic estimate and therefore strengthens the resulting
error lower bound.

However, the initial phase must be sufficiently long to ensure that
the adversarial problem belongs to the prescribed class $\Gaps_4$.
Its expected gaps must remain large enough after accounting for
fluctuations of the cumulative rewards.

We therefore seek the shortest initial phase, up to universal constant
factors, that guarantees these gap constraints uniformly over arms
satisfying the selection bounds. The choice below uses a convenient
factor $2$ to provide additional gap margin. Together with the phase
bound below, it keeps $\timeHorizon_i\gap_i^2/i$ bounded by a universal
constant times $\timeHorizon/\complexityBoth^{\basePro}$.

Set
\[
\dividefac_i\triangleq\min\left\{1,\frac{2\gap_1}{\gap_i}\right\},
\qquad 0<\dividefac_i\leq1.
\]
The initial phase $\earlyPhase_i$ is the prefix of rounds
$t=1,\ldots,\timeHorizon_i$, where
\[
\timeHorizon_i\triangleq\left\lceil\timeHorizon\dividefac_i\right\rceil
=\min\left\{\timeHorizon,
\left\lceil\frac{2\timeHorizon\gap_1}{\gap_i}\right\rceil\right\}.
\]
The equality uses the fact that $\timeHorizon$ is an integer;
in particular, $1\leq\timeHorizon_i\leq\timeHorizon$.

The theorem's horizon condition implies
\[
\timeHorizon\gap_1^2\geq128\log(32\nArms)\geq1.
\]
Writing $r=\timeHorizon\gap_1/\gap_i$, since
$\gap_1\gap_i\leq1/64$, we have
\[
r=\frac{\timeHorizon\gap_1^2}{\gap_1\gap_i}\geq1,
\qquad
\frac{\gap_i}{\timeHorizon\gap_1}
=\frac{\gap_1\gap_i}{\timeHorizon\gap_1^2}
\leq\frac{1}{64}\leq\frac18.
\]
Thus $\gap_i/\timeHorizon\leq\gap_1/8$, and the capped ceiling gives
\[
\timeHorizon_i\leq\lceil2r\rceil\leq2r+1\leq3r
=\frac{3\timeHorizon\gap_1}{\gap_i}.
\]

 
Since $\gap_1=\gap_2>0$, we have $1/\gap_1<2/\gap_2$.
Thus $i\geq2$ for every choice of maximizing index.

	Now we want to identify an arm with a gap of the same order as $\gap_i$ that the 
	learner will not pull
	very much in the base problem \basePro{}
	during the period~$\earlyPhase_i$. This arm will be the one 'attacked' by the adversary to confuse the learner as the learner is not pulling it very much.


If $i=2$, choose $\refarm=2$. Since at most one pull of arm $2$
is made in each round, $\pullsNumber_2(n_i)\leq n_i$ almost surely,
and hence
\[
\Exp_{\basePro}\left[\pullsNumber_2(n_i)\right]
\leq n_i=\frac{2n_i}{i}.
\]
Under the phase definition above, $i=2$ gives $\dividefac_i=1$
and $n_i=\timeHorizon$.
The second phase is then empty, and the \advPro{} construction below
coincides with \basePro{} throughout the horizon. This empty second
phase is allowed by the construction.

If $i\geq3$, consider the arms with indices in
$[\lceil i/2\rceil:i]$. There are $\lfloor i/2\rfloor+1>i/2$
such arms. Since
\[
\sum_{k\in\setArms}\Exp_{\basePro}\left[\pullsNumber_k(n_i)\right]=n_i,
\]
they cannot all have expected pull counts strictly greater than $2n_i/i$.
Choose any $\refarm\in[\lceil i/2\rceil:i]$ whose expected pull count
is at most $2n_i/i$. In this case $\refarm\geq\lceil i/2\rceil\geq2$.

 By construction we have $\gap_{\refarm}\leq\gap_i$ and
	%
	\begin{equation}\label{eq:notSoMuchPulls}
	\Exp_{\basePro}\left[\pullsNumber_{\refarm}(n_i) \right]\leq \frac{2n_i}{i}\cdot
	\end{equation}
	%    

In both cases, $2\leq\refarm\leq i$ and $\refarm\geq i/2$.
The maximizing property of $i$ gives
$i/\gap_i\geq\refarm/\gap_{\refarm}$, so
 \begin{equation}\label{eq:zhihan}
     \gap_i\leq\frac{i}{\refarm}\gap_{\refarm}\leq2\gap_{\refarm}.
 \end{equation}
Equality is permitted in this gap comparison. The resulting bound
$\gap_{\refarm}/\gap_i\geq1/2$ is sufficient for the expected-gap
lower bound below.

	
	
	Now we construct an i.i.d.\,stochastic problem, denoted \stoPro{}, where
	the arm distributions are the same as in the \basePro{}
	problem, except for arm $\refarm$, $\mu^{\stoPro}_{k}\triangleq \mu^{\basePro}_{k} $ 
	for all $k\neq \refarm$.
	%    We distinguish two cases: if the  $\refarm=1$ then
	We set in  \stoPro{},         
	$\mu^{\stoPro}_{\refarm} \triangleq 1/2 + \gap_1/4$.

The best arm in the constructed \stoPro{} problem is
$\bestArmSto\triangleq\refarm$. Its gaps, indexed by the original arm
labels, satisfy
\[
\gap_1^{\stoPro}=\gap_{\refarm}^{\stoPro}=\frac{\gap_1}{4},
\]
and, for every $j\in\setArms\setminus\{1,\refarm\}$,
\[
\gap_j\leq\gap_j^{\stoPro}
=\gap_j+\frac{\gap_1}{4}\leq\frac54\gap_j.
\]
Thus label $1$ has gap $\gap_1/4$, the selected label $\refarm$ is
the exceptional label with the same gap, and every remaining label
satisfies the ordinary bounds $[\gap_j/4,4\gap_j]$.
Consequently, $\stoPro\in\Gaps_4$.

We next compare the complexity of the constructed stochastic problem
with that of \basePro{}. Since the construction changes the ordering
of the arms, this comparison must use the sorted stochastic gaps.
Let $h_1\leq\cdots\leq h_\nArms$ be the sorted values of the labelled
gaps $\gap_j^{\stoPro}$, and define the local quantity
\[
Q\triangleq\max_{j\in\setArms}\frac{j}{\gap_j},
\qquad \complexityBoth^{\basePro}=\frac{Q}{\gap_1},
\qquad Q\geq\frac{2}{\gap_1}.
\]
The sorted stochastic gaps consist of two copies of $\gap_1/4$,
followed by the ordered list $\gap_j+\gap_1/4$ for
$j\in[2:\nArms]\setminus\{\refarm\}$.
Deleting the entry for $\refarm$ from the ordered list with labels
$2,\ldots,\nArms$ moves a remaining position by at most one.
This comparison remains valid when gaps tie. Hence
\[
h_1=h_2=\frac{\gap_1}{4},
\qquad
\gap_{r-1}+\frac{\gap_1}{4}\leq h_r
\leq\gap_r+\frac{\gap_1}{4}
\quad(3\leq r\leq\nArms).
\]
When $\nArms=2$, there are no ranks $r\geq3$; only the first two
ranks are used below.

For the upper bound, the first two ranks satisfy
\[
\max_{r\in\{1,2\}}\frac{r}{h_r}
=\frac{8}{\gap_1}\leq4Q.
\]
For each rank $r\geq3$, we have
\[
\frac{r}{h_r}\leq\frac{r}{\gap_{r-1}}
=\frac{r}{r-1}\frac{r-1}{\gap_{r-1}}
\leq\frac32Q\leq4Q.
\]
Thus, bounding every rank in $\setArms$ gives
\[
\complexityBoth^{\stoPro}
=\frac{4}{\gap_1}\max_{r\in\setArms}\frac{r}{h_r}
\leq\frac{4}{\gap_1}(4Q)
=16\complexityBoth^{\basePro}.
\]
For the lower bound, $h_r\leq2\gap_r$ for every rank $r$:
for $r=1,2$ this follows from $h_r=\gap_1/4$ and
$\gap_1=\gap_2$, while for $r\geq3$ it follows from
$h_r\leq\gap_r+\gap_1/4$ and $\gap_1\leq\gap_r$.
Consequently,
\[
\complexityBoth^{\stoPro}
\geq\frac{4}{\gap_1}\max_{r\in\setArms}\frac{r}{2\gap_r}
=2\complexityBoth^{\basePro}
\geq\frac12\complexityBoth^{\basePro}.
\]
Combining these bounds, we retain the weaker lower bound in
\begin{equation}\label{eq:newHBOBSTO}
  \complexityBoth^{\basePro} /2
	\leq
        \complexityBoth^{\stoPro} 
		\leq
 16 \complexityBoth^{\basePro} 
\end{equation}

Writing $\complexitySR^{\basePro}$ and $\complexitySR^{\stoPro}$
for the respective \SR{} complexities, the same sorted-gap comparison
and the stronger lower bound above give
\[
\begin{aligned}
\complexitySR^{\stoPro}
&\leq\max\left\{\frac{32}{\gap_1^2},
\frac32\complexitySR^{\basePro}\right\}
\leq16\complexitySR^{\basePro},\\
\frac{\complexityBoth^{\stoPro}}{\complexitySR^{\stoPro}}
&\geq\frac18\frac{\complexityBoth^{\basePro}}{\complexitySR^{\basePro}}.
\end{aligned}
\]
For the square-root profile $\gap_r=\gap_1\sqrt{r/2}$,
$r\in[2:\nArms]$, the ratio for \stoPro{} is at least
$\frac18\sqrt{\nArms/2}$, for every eligible choice of arm~$\refarm$.

	The expectation and probability with respect to the
	learner and the samples of this problem are denoted
	$\Exp_{\stoPro}$ and $\Pro_{\stoPro}$.
	

	
	The second bandit problem is the adversarial one, denoted \advPro{}.
Before interaction with the learner, it samples a table
$\gainVector\in\{0,1\}^{\nArms\times\timeHorizon}$.
For any sampled table with several maximizing arms, let
$\bestArm_{\gainVector}$ be the smallest maximizing index.
On $\eventc$ the maximizer is unique, so this convention does
not affect the extracted instance.
The Bernoulli entries are independent across arms and rounds,
and independent of the learner's internal randomness.
The table-generating distribution is fixed once the learner,
horizon, reference gaps, and the choices of $i$ and $\refarm$ are fixed;
it does not adapt to the learner's realized actions.
For each $t\in[\timeHorizon]$ and arm $k$, the $(k,t)$ entry is
Bernoulli with mean $\mean_k^{\advPro}(t)$.
	For all $t$ and $k\neq\refarm$, \advPro{} follows the \basePro{} 
	problem:  $\mu_{k}^{\advPro}(t)\triangleq\mu_{k}^{\basePro}$.
	For arm $\refarm$,
	until the end of phase $\earlyPhase_i$, for all $t$ with $1\leq 
	t\leq n_i$, $\mu^{\advPro}_{\refarm}(t)\triangleq\mu_{\refarm}^{\basePro}$ 
	and then a switch happens, for $n_i< t\leq n$, arm $\refarm$ 
	possesses  the same distributions as in the \stoPro{} problem,  
	$\mu^{\advPro}_{\refarm}(t)\triangleq \mu_{\refarm}^{\stoPro}$.
	The expectation and probability with respect to the
	learner and the samples of this problem are denoted
	$\Exp_{\advPro}$ and $\Pro_{\advPro}$.
	
	
Let us compute the expected cumulative gains for this fixed construction.
For each arm $k\in\setArms$, let
$\expectCumul_k\triangleq\sum_{t=1}^{\timeHorizon}\mean_k^{\advPro}(t)$.
Then
\[
\expectCumul_{\refarm}
=\timeHorizon_i\left(\frac12-\gap_{\refarm}\right)
+(\timeHorizon-\timeHorizon_i)\left(\frac12+\frac{\gap_1}{4}\right),
\qquad \expectCumul_1=\frac{\timeHorizon}{2},
\]
and, for $j\in\setArms\setminus\{1,\refarm\}$,
\[
\expectCumul_j=\timeHorizon\left(\frac12-\gap_j\right).
\]
Define the expected average gap
\[
b\triangleq\frac{\expectCumul_1-\expectCumul_{\refarm}}{\timeHorizon}
=\frac{\timeHorizon_i}{\timeHorizon}\gap_{\refarm}
-\left(1-\frac{\timeHorizon_i}{\timeHorizon}\right)\frac{\gap_1}{4}.
\]
If $\timeHorizon_i=\timeHorizon$, then $b=\gap_{\refarm}\geq\gap_1$.
This case also includes a full-horizon phase caused by ceiling rounding.
If $\timeHorizon_i<\timeHorizon$, its definition gives
$\timeHorizon_i\geq2\timeHorizon\gap_1/\gap_i$, so
Equation~\ref{eq:zhihan} yields
\[
b\geq\frac{2\gap_1\gap_{\refarm}}{\gap_i}-\frac{\gap_1}{4}
\geq\frac{3\gap_1}{4}.
\]
For the upper bound, in all cases,
\[
b\leq\frac{\timeHorizon_i}{\timeHorizon}\gap_{\refarm}
\leq\left(\frac{2\gap_1}{\gap_i}+\frac{1}{\timeHorizon}\right)\gap_i
=2\gap_1+\frac{\gap_i}{\timeHorizon}
\leq\frac{17\gap_1}{8}.
\]
Consequently,
\[
\frac{3\gap_1}{4}\leq b\leq\frac{17\gap_1}{8}.
\]

	%    
	Let us consider the event
	$\eventc =\left\{ \forall k\in\setArms,  |\cumulGain_k-\expectCumul_{k}| \leq  \timeHorizon \gap_1/16 \right\}
	$.
For a fixed arm $j$, $\cumulGain_j$ is the sum of all $\timeHorizon$
entries in its row of the reward table. These entries are independent
and lie in $[0,1]$, although their means may differ across rounds.
This is the full table sum, including rounds on which the learner
does not pull arm $j$. Two-sided Hoeffding gives
\begin{align*}
\Pro_{\advPro}\left(
|\cumulGain_j-\expectCumul_j|>\frac{\timeHorizon\gap_1}{16}\right)
&\leq2\exp\left(-\frac{2(\timeHorizon\gap_1/16)^2}{\timeHorizon}\right)\\
&=2\exp\left(-\frac{\timeHorizon\gap_1^2}{128}\right).
\end{align*}
Therefore, a union bound over arms yields
\begin{equation}\label{eq:goodEventBOB}
\Pro_{\advPro}(\eventc^c)
\leq2\nArms\exp\left(-\frac{\timeHorizon\gap_1^2}{128}\right)
\leq\frac{1}{16}.
\end{equation}
The last inequality uses the theorem hypothesis
\[
\nArms\exp\left(-\frac{\timeHorizon\gap_1^2}{128}\right)\leq\frac{1}{32}.
\]
Consequently, $\Pro_{\advPro}(\eventc)\geq15/16$.
%	\todov{here i put a power of 3 to be safe and i ttink it doesnt matter in the proof because we put it as an hypothesis of the theorem that $n$ is large enough}

On $\eventc$, each difference $\cumulGain_1-\cumulGain_j$
deviates from $\expectCumul_1-\expectCumul_j$ by at most
$\timeHorizon\gap_1/8$.
For the selected arm, the bounds on $b$ therefore give
\[
\frac{5\gap_1}{8}
\leq\frac{\cumulGain_1-\cumulGain_{\refarm}}{\timeHorizon}
\leq\frac{9\gap_1}{4}.
\]
For every $j\in\setArms\setminus\{1,\refarm\}$, we have
\[
\gap_j-\frac{\gap_1}{8}
\leq\frac{\cumulGain_1-\cumulGain_j}{\timeHorizon}
\leq\gap_j+\frac{\gap_1}{8}.
\]
Since $\gap_j\geq\gap_1$, this implies
\[
\frac78\gap_j
\leq\frac{\cumulGain_1-\cumulGain_j}{\timeHorizon}
\leq\frac98\gap_j.
\]
All these differences are strictly positive, so arm $1$ is the
unique best arm on $\eventc$. Thus the displayed differences are
the gaps $\gap_j^{\gainVector}$ for $j\neq1$.
The best arm's own gap is
\[
\gap_1^{\gainVector}
=\min_{j\in\setArms\setminus\{1\}}
\frac{\cumulGain_1-\cumulGain_j}{\timeHorizon}.
\]
The lower bounds for all other arms and the upper bound for the
selected arm give
\[
\frac{5\gap_1}{8}\leq\gap_1^{\gainVector}\leq\frac{9\gap_1}{4}.
\]
The bounds for labels $1$ and $\refarm$ lie in
$[\gap_1/4,4\gap_1]$, and the bounds for every remaining label $j$
lie in $[\gap_j/4,4\gap_j]$.
Hence every table $\gainVector\in\eventc$ belongs to $\Gaps_4$,
with $\refarm$ as the exceptional label in the class definition.

	Note that therefore the only difference between the two problems 
	\stoPro{} and \advPro{} is for  arm $\refarm$ during phase $\earlyPhase_i$.
	
	
The stochastic problem has best arm $\refarm$, whereas arm $1$ is
uniquely best in the adversarial construction on $\eventc$.
We compare the distributions of the learner's observations to quantify
the difficulty of distinguishing these problems.

Let $P_A$ and $P_S$ denote the laws of the observed transcript $Z$
from Lemma~\ref{l:transcriptChangeMeas} under randomized \advPro{}
and \stoPro{}, respectively. These transcript laws are distinguished
from the full table-and-learner law $\Pro_{\advPro}$ used for $\eventc$,
which also involves unobserved table entries.
Both constructions have independent Bernoulli entries, independent
of the learner's internal randomness, with all parameters in
$[3/8,17/32]\subset(0,1)$. The same learner is run with the same
horizon, and the only differing arm-round pairs are arm $\refarm$
through the initial prefix of $\timeHorizon_i$ rounds.
The arm and prefix are fixed before interaction.
Apply the lemma with $A=\advPro$, $B=\stoPro$,
$\mathcal T=\{1,\ldots,\timeHorizon_i\}$, and
\[
u=\frac12-\gap_{\refarm},\qquad v=\frac12+\frac{\gap_1}{4}.
\]
It gives
\[
D(P_A\Vert P_S)
=\Exp_{\advPro}\left[\pullsNumber_{\refarm}(\timeHorizon_i)\right]
\operatorname{kl}(u,v).
\]
Through this fixed prefix, \advPro{} and \basePro{} have identical
reward laws and use the same learner. Their interaction laws agree
through $\timeHorizon_i$, including the learner's adaptive actions.
Consequently,
\[
\Exp_{\advPro}\left[\pullsNumber_{\refarm}(\timeHorizon_i)\right]
=\Exp_{\basePro}\left[\pullsNumber_{\refarm}(\timeHorizon_i)\right],
\]
and hence
\[
D(P_A\Vert P_S)
=\Exp_{\basePro}\left[\pullsNumber_{\refarm}(\timeHorizon_i)\right]
\operatorname{kl}(u,v).
\]

For $u,v\in(0,1)$, the inequality $\log z\leq z-1$ gives
\begin{align*}
\operatorname{kl}(u,v)
&\leq u\left(\frac{u}{v}-1\right)
+(1-u)\left(\frac{1-u}{1-v}-1\right)\\
&=\frac{(u-v)^2}{v(1-v)}.
\end{align*}
Here $\gap_1\leq\gap_{\refarm}\leq1/8$, so
\[
|u-v|\leq\frac54\gap_{\refarm},\qquad
v(1-v)=\frac14-\frac{\gap_1^2}{16}\geq\frac{255}{1024}.
\]
Thus
\[
\operatorname{kl}(u,v)
\leq\frac{1600}{255}\gap_{\refarm}^2\leq8\gap_{\refarm}^2.
\]
Using Equation~\ref{eq:notSoMuchPulls}, $\gap_{\refarm}\leq\gap_i$,
the phase-length bound, and
$\complexityBoth^{\basePro}=i/(\gap_1\gap_i)$, we obtain
\begin{equation}\label{eq:transcriptKLBOB}
\begin{aligned}
D(P_A\Vert P_S)
&\leq8\gap_{\refarm}^2\frac{2\timeHorizon_i}{i}\\
&\leq\frac{16\gap_i^2\timeHorizon_i}{i}\\
&\leq\frac{48\timeHorizon\gap_1\gap_i}{i}\\
&=\frac{48\timeHorizon}{\complexityBoth^{\basePro}}\\
&\leq\frac{768\timeHorizon}{\complexityBoth^{\stoPro}}.
\end{aligned}
\end{equation}
The last inequality uses Equation~\ref{eq:newHBOBSTO}:
$\complexityBoth^{\stoPro}\leq16\complexityBoth^{\basePro}$ implies
$1/\complexityBoth^{\basePro}\leq16/\complexityBoth^{\stoPro}$.
The chosen phase is long enough to satisfy the adversarial gap
constraints, while its length bound keeps this information cost at
most a universal constant times $\timeHorizon/\complexityBoth^{\stoPro}$.

Assume the stochastic-error premise of Theorem~\ref{th:lowBOB}, and set
\[
x=\frac{\timeHorizon}{\complexityBoth^{\stoPro}}>0,\qquad
p=P_A(\recomArm_{\timeHorizon}=1),\qquad
q=P_S(\recomArm_{\timeHorizon}=1).
\]
Since the unique best arm in \stoPro{} is $\refarm\neq1$,
\[
q\leq\ProErr_{\stoPro}(\timeHorizon)
\leq\frac{1}{64}\exp(-2048x).
\]
If $q=0$, common support in Lemma~\ref{l:transcriptChangeMeas} gives
$p=0$. Otherwise $0<q<1$, and its event consequence for
$W=\{\recomArm_{\timeHorizon}=1\}$, together with
Equation~\ref{eq:transcriptKLBOB}, gives
\[
p\log(1/q)\leq D(P_A\Vert P_S)+\log2\leq768x+\log2.
\]
The stochastic guarantee also gives
$\log(1/q)\geq2048x+\log64=2048x+6\log2$. Therefore,
\[
p\leq\frac{768x+\log2}{2048x+6\log2}\leq\frac38.
\]
The denominator is positive, $768=(3/8)2048$, and
$\log2\leq(3/8)6\log2$, which proves the last inequality.
The same conclusion holds when $q=0$.
This bounds the probability of recommending arm $1$; that arm need
not be best on tables outside $\eventc$.

Return now to the full table-and-learner probability law, and write
\[
\ProErr_{\advPro}(\timeHorizon)
=\Pro_{\advPro}\left(\recomArm_{\timeHorizon}\neq\bestArm_{\gainVector}\right),
\]
with the fixed table tie-breaking convention above.
On $\eventc$, arm $1$ is uniquely best, and
$\Pro_{\advPro}(\recomArm_{\timeHorizon}=1)=p$ because this is a
transcript event. Equation~\ref{eq:goodEventBOB} therefore yields
\begin{equation}\label{eq:randomizedErrorBOB}
\begin{aligned}
\ProErr_{\advPro}(\timeHorizon)
&\geq\Pro_{\advPro}\left(\{\recomArm_{\timeHorizon}\neq1\}\cap\eventc\right)\\
&\geq\Pro_{\advPro}(\recomArm_{\timeHorizon}\neq1)
-\Pro_{\advPro}(\eventc^c)\\
&=1-p-\Pro_{\advPro}(\eventc^c)\\
&\geq1-\frac38-\frac{1}{16}\\
&=\frac{9}{16}\geq\frac18.
\end{aligned}
\end{equation}
No independence between the recommendation and $\eventc$ is needed.

Equation~\ref{eq:randomizedErrorBOB} gives the lower bound needed
for the extraction of a deterministic table:
\[
\ProErr_{\advPro}(\timeHorizon)\geq1/8.
\]
For a fixed table $\gainVector$, write
\[
\ProErr_{\gainVector}(\timeHorizon)
\triangleq\ProErr_{\advPro(\gainVector)}(\timeHorizon),
\]
where the probability is over the learner's internal randomness.
The random \advPro{} model is a distribution over fixed tables,
independent of that randomness, so
\[
\ProErr_{\advPro}(\timeHorizon)
=\Exp_{\gainVector\sim\advPro}
\left[\ProErr_{\gainVector}(\timeHorizon)\right].
\]
Suppose for contradiction that $\ProErr_{\gainVector}(\timeHorizon)<1/16$
for every $\gainVector\in\eventc$. Splitting this expectation using
indicators and applying Equation~\ref{eq:goodEventBOB} gives
\begin{align*}
\ProErr_{\advPro}(\timeHorizon)
&=\Exp_{\gainVector\sim\advPro}\left[
\ProErr_{\gainVector}(\timeHorizon)\indicator{\gainVector\in\eventc}\right]
+\Exp_{\gainVector\sim\advPro}\left[
\ProErr_{\gainVector}(\timeHorizon)\indicator{\gainVector\notin\eventc}\right]\\
&\leq\frac{1}{16}\Pro_{\advPro}(\eventc)
+\Pro_{\advPro}(\eventc^c)\\
&=\frac{1}{16}+\frac{15}{16}\Pro_{\advPro}(\eventc^c)\\
&\leq\frac{1}{16}+\frac{15}{256}\\
&=\frac{31}{256}<\frac18.
\end{align*}
This contradicts the stated input $\ProErr_{\advPro}(\timeHorizon)\geq1/8$.
Hence there exists a deterministic table $\gainVector\in\eventc$ with
$\ProErr_{\advPro(\gainVector)}(\timeHorizon)\geq1/16$.
The gap calculation above shows that this table belongs to $\Gaps_4$.
	%
\end{proof}%
